\section{Lattice fields}\label{sec:uv-lattice}

The field at a point will be represented by observables on finitely many nearby links.
There are two choices to make: which classical field an observable represents, and how to change its coefficients as the lattice spacing decreases.
This section specifies these choices without yet integrating out any scales.

We begin with the exact symmetries of the Wilson measure in Subsection~\ref{subsec:uv-law}.
Subsection~\ref{subsec:uv-labels} describes classical local fields by their Taylor coefficients and chooses finite bases of polynomial labels.
Subsection~\ref{subsec:uv-local} states how a polynomial in link matrices must reproduce those labels.
Subsection~\ref{subsec:uv-renormalized} then gives the reference-scale normalization of the lattice insertions.
Section~\ref{sec:uv-activities} constructs the local polynomials and introduces the activity sums used in the integration argument; the coefficient transfers require the later renormalization steps.

We use the lattice, plaquettes and measure of Definitions~\ref{def:lattice}, \ref{def:plaquette} and~\ref{def:wilson-law}.
The gauge group is as in Definition~\ref{def:gauge}; simplicity and connectedness are not needed in this section.

\subsection{The Wilson measure}\label{subsec:uv-law}

The action depends on the transport around each elementary square, or plaquette, as explained in Definition~\ref{def:plaquette}.
Conjugating that transport or reversing its direction leaves the action unchanged.
This observation gives exact symmetries of the probability measure, before any limit is taken.

Write~$U(w)$ for the ordered transport along a walk~$w$ from~$x$ to~$z$, using the link convention of Definition~\ref{def:lattice}.
The gauge action of Definition~\ref{def:gauge-transform} gives
\begin{equation}\label{eq:uv-gauge}
 U^\gamma(w)=\gamma(x)U(w)\gamma(z)^{-1}\,.
\end{equation}
The factors at intermediate vertices cancel, for forward and backward steps alike.
For the estimates below, abbreviate the single-plaquette action and define its associated distance from the identity by
\begin{equation}\label{eq:uv-chord}
 s(u)=1-\frac{1}{m}\operatorname{Re}\tr u,
 \qquad \chi(u)=\frac{\|u-1\|_{\mathrm F}}{\sqrt m}\qquad(u\in G),
\end{equation}
where~$\|\cdot\|_{\mathrm F}$ is the Frobenius norm.
Unitarity gives
\begin{equation}\label{eq:uv-chord-square}
 \chi(u)^2=2s(u),\qquad 0\le s(u)\le2,\qquad
 \chi(uv)\le\chi(u)+\chi(v)\,.
\end{equation}
Indeed, $\|u-1\|_{\mathrm F}^2=2m-2\operatorname{Re}\tr u$ and~$\|u-1\|_{\mathrm F}\le2\sqrt m$.
For the last inequality, write~$uv-1=(u-1)+u(v-1)$ and use invariance of the Frobenius norm under unitary multiplication.
Both~$s$ and~$\chi$ are continuous and invariant under conjugation and inversion.

The coordinate symmetries we use are the permutations and reversals of the four lattice axes.
They act on link configurations by transporting each link along its image.

\begin{definition}[Coordinate symmetries]\label{def:uv-symmetries}
Let~$H_4$ be the group of signed permutation matrices.
For~$A(y)=hy+z$, with~$h\in H_4$ and~$z\in\Lambda_\ell$, the pullback~$A^*U$ reads the transport along the image of each positive link.
A reversed image uses the inverse variable, with its initial vertex shifted accordingly.
The lattice and link conventions are those of Definition~\ref{def:lattice}.
\end{definition}

Write~$\mathcal U_\ell$ for the configuration space of Definition~\ref{def:lattice}, and~$B_b(\mathcal U_\ell;\C)$ for its bounded Borel complex functions.
The following proposition justifies ordinary expectations of local polynomials and the changes of variables used throughout the construction.

\begin{proposition}[The lattice law]\label{prop:uv-law}
For~$\ell\ge1$ and~$\beta\in\R$, the measure of Definition~\ref{def:wilson-law} satisfies
\begin{equation}\label{eq:uv-law}
 \mu_{\ell,\beta}(\mathcal U_\ell)=1,\qquad
 C(\mathcal U_\ell;\C)\cup B_b(\mathcal U_\ell;\C)
 \subseteq L^1(\mu_{\ell,\beta})\,.
\end{equation}
For every gauge transformation~$\gamma$ of Definition~\ref{def:gauge-transform}, coordinate symmetry~$A$ of Definition~\ref{def:uv-symmetries}, and~$F\in L^1(\mu_{\ell,\beta})$,
\begin{equation}\label{eq:uv-expectation-symmetry}
 \int F(U^\gamma)\,d\mu_{\ell,\beta}(U)
 =\int F(U)\,d\mu_{\ell,\beta}(U)
 =\int F(A^*U)\,d\mu_{\ell,\beta}(U)\,.
\end{equation}
\end{proposition}

\begin{proof}
We preserve the reference product measure and the energy separately.
Each transformation therefore preserves the Wilson measure.
Write~$\eta_\ell$ for the product Haar measure in~\eqref{eq:wilson-law}.

The closed group~$G\subset U(m)$ and its finite configuration space are compact.
There are six plaquettes per site, each with energy in~$[0,2]$.
Thus
\begin{equation}\label{eq:uv-density-bounds}
 0\le E_\ell\le12\ell^4,\qquad
 e^{-12|\beta|\ell^4}\le e^{-\beta E_\ell(U)}\le e^{12|\beta|\ell^4}\,.
\end{equation}
Integrating against probability Haar gives a finite positive normalizer.
A continuous function on the configuration space is bounded, so both classes in~\eqref{eq:uv-law} are integrable.

Haar invariance allows a fixed matrix to be multiplied on either side of a link, and allows its orientation to be reversed.
For~$R_g(u)=ug$, $\mathrm{inv}(u)=u^{-1}$ and every integrable~$f:G\to\C$, the identities are
\begin{equation}\label{eq:uv-haar-changes}
 (R_g)_*h=h,\qquad \mathrm{inv}_*h=h,\qquad
 \int_G f(aub)\,h(du)=\int_G f(u)\,h(du)\quad(a,b\in G)\,.
\end{equation}
Indeed, $(R_g)_*h$ is left invariant with mass one, so Haar uniqueness identifies it with~$h$.
The image under inversion is then left invariant by right invariance of~$h$, again with mass one.
Applying these identities independently to each link gives, for the gauge map~$\Gamma_\gamma:U\mapsto U^\gamma$,
\begin{equation}\label{eq:uv-product-changes}
 (\Gamma_\gamma)_*\eta_\ell=\eta_\ell,\qquad
 (A^*)_*\eta_\ell=\eta_\ell\,.
\end{equation}
For the second identity, the coordinate symmetry permutes the positive links and inverts those whose image orientation is negative.

It remains to check the energy.
By~\eqref{eq:uv-gauge}, a gauge transformation conjugates a plaquette transport at its initial vertex.
A coordinate symmetry sends an elementary square to another one, up to a change of initial corner and orientation.
To check the counting, write~$he_\mu=\sigma_\mu e_{\pi(\mu)}$, with~$\sigma_\mu\in\{-1,1\}$.
The positive-axis base of the image square is
\begin{equation*}
 hy+z-\mathbf1_{\{\sigma_\mu=-1\}}e_{\pi(\mu)}
      -\mathbf1_{\{\sigma_\nu=-1\}}e_{\pi(\nu)},
\end{equation*}
and its axis pair is the increasing rearrangement of~$\{\pi(\mu),\pi(\nu)\}$.
This permutes the face indices even when~$\ell=1$ or~$2$.
Since conjugation and inversion preserve~$s$, we have
\begin{equation}\label{eq:uv-energy-invariance}
 E_\ell(U^\gamma)=E_\ell(U),\qquad E_\ell(A^*U)=E_\ell(U)\,.
\end{equation}
For either~$T=\Gamma_\gamma$ or~$T=A^*$, equations~\eqref{eq:uv-product-changes} and~\eqref{eq:uv-energy-invariance} now give
\begin{equation*}
 \int F(TU)\,d\mu_{\ell,\beta}(U)
 =\frac{1}{Z_{\ell,\beta}}\int F(TU)e^{-\beta E_\ell(TU)}\,d\eta_\ell(U)
 =\int F(U)\,d\mu_{\ell,\beta}(U)\,.
\end{equation*}
This proves~\eqref{eq:uv-expectation-symmetry} for every integrable~$F$.
\end{proof}

\subsection{Classical labels}\label{subsec:uv-labels}

We now choose concrete classical quantities for the lattice observable labels of Definition~\ref{def:observables}.
A connection near the origin is described in radial gauge: its component in the radial direction vanishes.
Its Taylor coefficients then satisfy a linear constraint, degree by degree.
The constant coefficient vanishes: a constant one-form with zero contraction against every radial vector must be zero.
A polynomial in finitely many such coefficients describes a classical local observable; invariance under simultaneous conjugation removes the choice of gauge at the origin.

\begin{definition}[Classical labels]\label{def:uv-labels}
For~$\g$ of Definition~\ref{def:gauge} and~$j\ge0$, let~$\mathcal A_j$ be the real space of homogeneous polynomial one-forms of degree~$j+1$, with values in~$\g$, satisfying
\begin{equation}\label{eq:uv-radial}
 \sum_{\mu=1}^4 z^\mu A_{j,\mu}(z)=0\qquad(z\in\R^4)\,.
\end{equation}
The list of Taylor coefficients, called a real jet, is a sequence~$J=(A_j)_{j\ge0}$ with~$A_j\in\mathcal A_j$.
Choose real coordinates on each~$\mathcal A_j$, assigning weight~$j+2$ to each coordinate.
Let~$V_d$ be the space of complex polynomials in these coordinates, of total weight at most~$d$, invariant under simultaneous conjugation by~$G$.
\end{definition}

The coefficients in Definition~\ref{def:uv-labels} are Taylor coefficients, with the factorial normalization, rather than derivative values.
Equivalently, they are diagonal evaluations of symmetric~$(j+1)$-linear maps.
The weight~$j+2$ records the scaling of a link integral: the polynomial one-form contributes~$j+1$ powers of length and integration contributes one more.
Weights add under multiplication.
We also count the number of factors in a monomial, its ordinary polynomial degree; the two gradings will serve different purposes.

An orthogonal matrix~$h$ acts on these one-forms and on their invariant polynomials by
\begin{equation}\label{eq:uv-label-action}
 (T_hA_j)_\mu(z)=\sum_\nu h_{\mu\nu}A_{j,\nu}(h^{-1}z),
 \qquad (h\cdot v)(J)=v(T_{h^{-1}}J)\,.
\end{equation}
This action preserves~\eqref{eq:uv-radial} and commutes with conjugation by~$G$.
Coefficient conjugation is denoted by~$v\mapsto\bar v$; a real polynomial has~$\bar v=v$.
Let~$V_d^{w,f}$ be the subspace homogeneous of weight~$w$ and factor number~$f$.

The weight bound leaves only finitely many coordinates and monomials.
We can therefore choose one finite basis at each degree, retaining the basis elements already chosen at lower degrees.
Real homogeneous bases make conjugation and the two gradings explicit.

\begin{lemma}[Finite label spaces]\label{lem:uv-labels}
For the spaces of Definition~\ref{def:uv-labels}, evaluation on real jets is injective, and
\begin{equation}\label{eq:uv-label-grading}
 \dim_\C V_d<\infty,\qquad 1\in V_d\subseteq V_{d+1},\qquad
 V_d=\bigoplus_{\substack{0\le w\le d\\0\le f\le w}}V_d^{w,f}\,.
\end{equation}
Each summand is preserved by~\eqref{eq:uv-label-action} and coefficient conjugation.
There are nested real bases, homogeneous in both gradings and containing~$1$.
\end{lemma}

\begin{proof}
The weight bound reduces the problem to finitely many monomials; we then choose bases separately in their real homogeneous components.

Only coordinates with~$j+2\le d$ can occur in a monomial of weight at most~$d$.
There are finitely many such coordinates, and the exponent of each is bounded by~$d$.
Thus finitely many monomials span the polynomials under consideration, and their invariant subspace is finite dimensional.
A complex polynomial vanishing on every real jet vanishes identically, by the one-variable polynomial identity theorem applied successively in these finitely many real coordinates.

Conjugation acts linearly within each~$\mathcal A_j$ and therefore preserves both grades.
Writing~$v=\sum_{w,f}v^{w,f}$, comparison of the homogeneous terms gives
\begin{equation}\label{eq:uv-homogeneous-invariance}
 v(\Ad(g)J)=v(J) \Longrightarrow\quad
 v^{w,f}(\Ad(g)J)=v^{w,f}(J)\qquad(g\in G)\,.
\end{equation}
Thus the homogeneous parts of an invariant polynomial are themselves invariant.
Every coordinate has weight at least two; in particular the factor number cannot exceed the weight, giving the index range in~\eqref{eq:uv-label-grading}.
The same linearity argument applies to~\eqref{eq:uv-label-action}, and these real linear changes commute with coefficient conjugation.
Each invariant homogeneous component is the complex span of its real polynomials, by decomposition into real and imaginary coefficient parts.
Choose a real basis in each component, beginning with~$\{1\}$ for the constant component, and retain these choices as~$d$ increases.
\end{proof}

Fix the bases of Lemma~\ref{lem:uv-labels}, writing~$b_{d,i}$ for their elements and~$(w_{d,i},f_{d,i})$ for their two grades, with~$b_{d,0}=1$.
We identify each polynomial with its coefficient vector when applying a linear map to~$V_d$.
The continuum spaces~$F_d$ of Definition~\ref{def:hierarchy} will arise only after a limiting correlation table and its null relations have been obtained.
The spaces~$V_d$ still contain all classical labels, before these relations are imposed.

\subsection{Local representatives}\label{subsec:uv-local}

A polynomial in Taylor coefficients is not yet a lattice observable.
We seek a polynomial in nearby link matrices whose Taylor expansion, when tested on a classical connection, recovers the chosen label.
This is the purpose of a local section.
Its existence is a separate construction; here we state its normalization precisely.

On~$\Z^4$, the gauge action of Definition~\ref{def:gauge-transform} and the coordinate pullbacks of Definition~\ref{def:uv-symmetries} use the same formulas as on a torus.
Polynomial functions on a finite set of links can then be evaluated on any periodic lattice.

\begin{definition}[Local polynomial functions]\label{def:uv-polynomials}
For a finite set~$\Sigma\subset\Z^4$, let~$\mathcal P(\Sigma)$ be the polynomial functions of matrix entries and conjugate entries of links with both endpoints in~$\Sigma$.
At~$o\in\Lambda_\ell$, evaluate such a polynomial by substituting
\begin{equation}\label{eq:uv-placement}
 U(x,\mu)\longmapsto U(o+[x]_\ell,\mu)\,.
\end{equation}
The link conventions are those of Definition~\ref{def:lattice}.
\end{definition}

This implements the finite-site locality requirement~\ref{reg:4} of Definition~\ref{def:regular}.
The substitution is meaningful even when distinct cover links become the same torus link.

To test the classical normalization, truncate a real jet of Definition~\ref{def:uv-labels} to weights at most~$d$ and set~$A^J_\mu=\sum_{j+2\le d}A_{j,\mu}$.
We obtain the test links by solving a matrix equation, and then use closedness of~$G$ to see that the solution remains in the group.
For a spacing parameter~$\varepsilon\in\R$, put
\begin{equation}\label{eq:uv-link-ode}
 \begin{aligned}
 B(t)&=\varepsilon A_\mu^J\bigl(\varepsilon(x+te_\mu)\bigr), & u'(t)&=u(t)B(t),\\
 u(0)&=1, & U_{J,\varepsilon}(x,\mu)&=u(1)\,.
 \end{aligned}
\end{equation}
Successive integration gives the ordered series
\begin{equation}\label{eq:uv-transport-series}
 \begin{aligned}
 u(t)&=1+\sum_{r\ge1}\int_{0<s_1<\cdots<s_r<t}
          B(s_1)\cdots B(s_r)\,ds_1\cdots ds_r,\\
 \left\|\int_{0<s_1<\cdots<s_r<t}B(s_1)\cdots B(s_r)\,ds_1\cdots ds_r\right\|_{\mathrm F}
 &\le\frac{M^r}{r!}
 \end{aligned}
\end{equation}
for~$r\ge1$ and~$0\le t\le1$ when~$\|B(s)\|_{\mathrm F}\le M$.
For a fixed real jet, $B$ is polynomial in~$\varepsilon$ and the same factorial majorant holds on a small complex~$\varepsilon$-neighborhood of zero.
The series therefore converges normally there, solves the equation, and gives an endpoint real analytic in~$\varepsilon$ near zero.
Iterating the integral equation for the difference of two solutions gives the same factorial bound and proves uniqueness.
For real~$\varepsilon$, step-function approximations to~$B$ have solutions that are products of matrices~$\exp(\Delta t\,B(t_j))\in G$.
The same series bound proves convergence of these products to~$u$; since~$G$ is closed, the links in~\eqref{eq:uv-link-ode} are group valued.

The section must preserve the exact lattice symmetries and the constant observable.
Its classical normalization requires equality of Taylor coefficients through the degree being represented.
These are the conditions in the next definition.

\begin{definition}[Normalized local section]\label{def:uv-sections}
A normalized local section in degree~$d$ is a finite set~$\Sigma_d\subset\Z^4$ and a complex linear map~$P_d:V_d\to\mathcal P(\Sigma_d)$, with polynomial functions as in Definition~\ref{def:uv-polynomials}, such that
\begin{equation}\label{eq:uv-section-support}
 0\in\Sigma_d,\qquad h\Sigma_d=\Sigma_d\quad(h\in H_4),
\end{equation}
and, for all~$v\in V_d$, link configurations~$U$ on~$\Z^4$, $\gamma:\Z^4\to G$ and~$h\in H_4$,
\begin{equation}\label{eq:uv-section-symmetry}
 \begin{aligned}
 P_d[v](U^\gamma)&=P_d[v](U), & P_d[1]&=1,\\
 P_d[\bar v](U)&=\overline{P_d[v](U)}, &
 P_d[v](h^*U)&=P_d[h\cdot v](U)\,.
 \end{aligned}
\end{equation}
For every real jet~$J$ and~$0\le k\le d$, it also satisfies
\begin{equation}\label{eq:uv-section-taylor}
 \left.\frac{d^k}{d\varepsilon^k}P_d[v](U_{J,\varepsilon})\right|_{\varepsilon=0}
 =\left.\frac{d^k}{d\varepsilon^k}v(\operatorname{sc}_\varepsilon J)\right|_{\varepsilon=0},
 \qquad (\operatorname{sc}_\varepsilon J)_j=\varepsilon^{j+2}A_j\,.
\end{equation}
Here~$V_d$, its real jets and its action are as in Definition~\ref{def:uv-labels} and~\eqref{eq:uv-label-action}; $H_4$ and its pullbacks are from Definition~\ref{def:uv-symmetries}, and the test links are~\eqref{eq:uv-link-ode}.
\end{definition}

The Taylor identity imposes no equality at nonzero spacing.
No compatibility of~$P_d$ and~$P_{d+1}$ is required.
Proposition~\ref{prop:uv-sections} constructs these sections and proves the derivative estimates needed for the later blocking argument.

\subsection{Renormalized insertions}\label{subsec:uv-renormalized}

As shorter scales are integrated out, the coefficients of an observable change and different classical labels can mix.
We normalize the observable by its coefficients at one reference level~$r(d)$, chosen independently of cutoff and volume.
The inverse of the coefficient transfer then determines which fine-lattice polynomial represents a prescribed reference label.
This describes a particular choice of the insertions in Definition~\ref{def:observables}, not an additional condition on that definition.

Use~$a_N,N,p$ from Definition~\ref{def:family}, the basis from Lemma~\ref{lem:uv-labels}, and a local section from Definition~\ref{def:uv-sections}.
Let~$\theta_N>0$ be the threshold that will distinguish small plaquette distances~$\chi$ in the later estimates.
Define the linear factor scaling by~$S_t b_{d,i}=t^{f_{d,i}}b_{d,i}$.
For an invertible complex linear map~$\tau_{N,p,d}:V_d\to V_d$, the candidate insertion is
\begin{equation}\label{eq:uv-insertion}
 I_{N,p}\bigl((d,v),o;U\bigr)
 =\frac{1}{a_N^d}P_d\!\left[S_{\theta_N^{-1}}\tau_{N,p,d}^{-1}v\right](o;U),
\end{equation}
with placement~\eqref{eq:uv-placement}.
Thus a label with~$f$ factors is multiplied by~$\frac{1}{\theta_N^f}$; the prefactor~$\frac{1}{a_N^d}$ supplies the physical degree normalization.

For~$N\ge r(d)$, let~$x_j:V_d\to V_d$ be complex linear maps from coefficients at the finest level~$N$ to coefficients at level~$j$, for~$r(d)\le j\le N$.
The conditions that make~$\tau_{N,p,d}$ a reference transfer are
\begin{equation}\label{eq:uv-reference-coordinates}
 x_N(v)=v,\qquad x_{r(d)}(v)=\tau_{N,p,d}v,\qquad
 x_j(a_N^d1)=a_j^d1\quad(r(d)\le j\le N)\,.
\end{equation}
These equalities fix how coefficients are normalized between levels; they assert no pointwise identity between observables on different lattices.
Reality and signed-permutation covariance require
\begin{equation}\label{eq:uv-transfer-symmetry}
 \tau_{N,p,d}\bar v=\overline{\tau_{N,p,d}v},\qquad
 \tau_{N,p,d}(h\cdot v)=h\cdot(\tau_{N,p,d}v)\quad(h\in H_4)\,.
\end{equation}

The locality of the resulting insertion follows from its polynomial formula, even on a torus smaller than the chosen set of sites.
The unit normalization is different: it uses the reference-transfer conditions and therefore applies only beyond the reference level.
The following proposition separates these two conclusions.

\begin{proposition}[Locality of the insertions]\label{prop:uv-insertions}
Let~$P_d$ satisfy Definition~\ref{def:uv-sections}, $\theta_N>0$ and~$\tau_{N,p,d}\in\mathrm{GL}_\C(V_d)$.
For all~$N,p$, the map~$v\mapsto I_{N,p}\bigl((d,v),o;\cdot\bigr)$ in~\eqref{eq:uv-insertion} is complex linear and satisfies~\ref{reg:2}--\ref{reg:4} of Definition~\ref{def:regular}, with finite set~$\Sigma_d$.
Moreover,
\begin{equation}\label{eq:uv-insertion-translation}
 I_{N,p}\bigl((d,v),o;U(\cdot+z)\bigr)=I_{N,p}\bigl((d,v),o+z;U\bigr)\,.
\end{equation}
If~$N\ge r(d)$ and~\eqref{eq:uv-reference-coordinates}--\eqref{eq:uv-transfer-symmetry} hold, then
\begin{equation}\label{eq:uv-insertion-covariance}
 \begin{aligned}
 I_{N,p}\bigl((d,\bar v),o;U\bigr)&=\overline{I_{N,p}\bigl((d,v),o;U\bigr)},\\
 I_{N,p}\bigl((d,h\cdot v),A(o);U\bigr)&=I_{N,p}\bigl((d,v),o;A^*U\bigr),\\
 I_{N,p}\bigl((d,a_{r(d)}^d1),o;U\bigr)&=1,\qquad
 x_{r(d)}(\tau_{N,p,d}^{-1}v)=v
 \end{aligned}
\end{equation}
for~$A(o)=ho+z$, $h\in H_4$, as in Definition~\ref{def:uv-symmetries}.
\end{proposition}

\begin{proof}
Polynomial dependence gives continuity and finite-site locality; the endpoint condition then identifies which label represents the constant observable.

The coefficient changes in~\eqref{eq:uv-insertion} are complex linear.
A polynomial in link entries and their conjugates is continuous, and the section is gauge invariant by~\eqref{eq:uv-section-symmetry}; this proves~\ref{reg:2} and~\ref{reg:3}.
Placement~\eqref{eq:uv-placement} reads only links with both endpoints in~$o+[\Sigma_d]_{L_pD_N}$, proving~\ref{reg:4}.
Identifying distinct cover vertices does not enlarge that set.
Translating the placement proves~\eqref{eq:uv-insertion-translation}.

Factor scaling commutes with coefficient conjugation and~$H_4$, since both actions preserve factor number by Lemma~\ref{lem:uv-labels}.
Together with~\eqref{eq:uv-transfer-symmetry} and~\eqref{eq:uv-section-symmetry}, this gives the first two identities in~\eqref{eq:uv-insertion-covariance}.
For the unit, the endpoint conditions~\eqref{eq:uv-reference-coordinates} give
\begin{equation}\label{eq:uv-unit-calculation}
 \begin{aligned}
 \tau_{N,p,d}^{-1}(a_{r(d)}^d1)&=a_N^d1,\\
 I_{N,p}\bigl((d,a_{r(d)}^d1),o;U\bigr)
 &=\frac{1}{a_N^d}P_d\bigl[S_{\theta_N^{-1}}(a_N^d1)\bigr](o;U)
 =\frac{a_N^d}{a_N^d}P_d[1](o;U)=1\,.
 \end{aligned}
\end{equation}
The remaining identity follows by applying the reference endpoint condition to~$\tau_{N,p,d}^{-1}v$.
\end{proof}

Choose the lattice observable labels in Definition~\ref{def:observables} to be~$\mathcal T=\bigsqcup_dV_d$.
Under the hypotheses of Proposition~\ref{prop:uv-insertions}, the requirements~\ref{reg:2}--\ref{reg:4} are satisfied at every cutoff; the reference unit is~$a_{r(d)}^d1$.
Proposition~\ref{prop:uv-law} makes every lattice moment in~\eqref{eq:smear} an integral of a continuous, hence integrable, product of observables.
The running-coupling condition~\ref{reg:1} and convergence of these moments require the later ultraviolet arguments.

\gap{The invertible coefficient transfers satisfying~\eqref{eq:uv-reference-coordinates}--\eqref{eq:uv-transfer-symmetry} must still be constructed for the same Wilson family. Proposition~\ref{prop:uv-insertions} assumes them. The local sections and their derivative bounds are supplied by Proposition~\ref{prop:uv-sections}; transferring their coefficients through the integration over scales remains a later ultraviolet proof.\sourcefootnote{Transfer conditions: \texttt{YM/Source/Insertion.lean}, lines 89--114; same-family application: \texttt{YM/FullVolume/Cross.lean}, lines 30--71}}
